\section{Operator-sufficient diffusion beliefs}
\label{sec:method}

\subsection{Problem and local belief}

Let $x\in\R^n$ be the target, with acquired data
$y_M=A_Mx+\epsilon_M$.  Action $a\in\mathcal A(M)$ reveals
\begin{equation}
    y_a=A_ax+\epsilon_a,\qquad
    \epsilon_a\sim\N(0,\Gamma_a).
    \label{eq:action-model}
\end{equation}
$A_a$ may be a complete complex multicoil phase-encode line in stacked real
coordinates, or selection of range/reflectance values in an angular LiDAR
tile.  $\Gamma_a$ is positive diagonal; a known whitening operator also
suffices.

At selected reverse-diffusion checkpoints, maintain $K$ trajectories
$x_t^k$ guided by $y_M$.  Their clean predictions are
$m_k=D_\theta(x_t^k,t)$.  A schedule-correct Tweedie coefficient $\kappa_t$
relates the local denoiser Jacobian to a conditional covariance surrogate.
Rather than form it in image space, we push $r_k$ candidate-space probes
through central denoiser differences.  A protected Nystr\"om factorization
produces aligned factors
\begin{equation}
 \begin{aligned}
 x\mid Z=k &\approx m_k+U_k\xi_k,\quad
 \xi_k\sim\N(0,I_{r_k}),\\[-1mm]
 \Pr(Z=k)&=w_k .
 \end{aligned}
    \label{eq:branch}
\end{equation}
The same coefficient $\xi_k$ must generate the state and every candidate
response.  This alignment, not merely matching marginal variances, is what
makes analytic conditioning possible.

For a complete candidate bank $B$, define
\begin{equation}
    \mu_k=Bm_k,\qquad V_k=BU_k.
    \label{eq:bank}
\end{equation}
Only products with $B$ and $B^\top$ are required: $B$ is full SENSE,
$(Bx)_c=\mathcal F(S_cx)$, for MRI and the identity before LiDAR tile
indexing.  Jacobian bias, finite differences, PSD projection, and imperfect
guidance make Eq.~\eqref{eq:branch} a surrogate, not an exact posterior.

\subsection{Information decomposition and bounds}

For action indices $P_a$, set
\begin{equation}
 \mu_{ka}=P_a\mu_k,\qquad
 \widetilde V_{ka}=\Gamma_a^{-1/2}P_aV_k.
\end{equation}
Under Eq.~\eqref{eq:branch},
$y_a\mid Z=k$ is Gaussian.  Information about the represented latent
$C=(Z,\xi_Z)$ decomposes as
\begin{align}
 \I(C;y_a)
 &=L_a+M_a, \label{eq:mi-decomposition}\\
 L_a
 &=\sum_k\frac{w_k}{2}
   \log\det(I+\widetilde V_{ka}^{\top}\widetilde V_{ka}),\\
 M_a&=\I(Z;y_a).
\end{align}
$L_a$ is analytic; only the categorical term is sampled from the small
predictive mixture.  Each component width comes from its local diffusion
tangent, rather than an isotropic kernel around endpoint particles.

Let $\bar\mu_a=\sum_kw_k\mu_{ka}$ and concatenate
\begin{equation}
 G_a=\left[
   \sqrt{w_k}P_aV_k,\;
   \sqrt{w_k}(\mu_{ka}-\bar\mu_a)
 \right]_{k=1}^K.
\end{equation}
The total-covariance identity gives $\Cov(P_aBx)=G_aG_a^\top$ and
\begin{equation}
 \boxed{\begin{aligned}
 L_a&\leq \I(C;y_a)\leq U_a,\\
 U_a&=\min\{U_a^{\rm G},\,L_a+H(w)\},\\
 U_a^{\rm G}&=\tfrac12\log\det
 (I+\widetilde G_a^\top\widetilde G_a).
 \end{aligned}}
 \label{eq:sandwich}
\end{equation}
where $\widetilde G_a=\Gamma_a^{-1/2}G_a$ and
$H(w)=-\sum_kw_k\log w_k$.  Thus $U_a<\max_bL_b$ safely removes $a$ for a
\emph{one-step, equal-cost} information objective.  A fixed candidate cap is
an engineering truncation and is recorded separately.

\subsection{Sufficient action coordinates}

Choose an orthonormal basis
\begin{equation}
    Q_a=\orth(\widetilde G_a),\qquad
    d_a=\rank(Q_a).
    \label{eq:q}
\end{equation}
Whiten and project an observation:
\begin{equation}
    s_a=Q_a^\top\Gamma_a^{-1/2}(y_a-\bar\mu_a).
    \label{eq:projected-observation}
\end{equation}
Coordinates orthogonal to $Q_a$ are component-independent standard Gaussian
noise and can be discarded.  In projected coordinates,
\begin{equation}
  s_a\mid k\sim\N(\delta_{ka},I+R_{ka}R_{ka}^\top),
\end{equation}
with
\begin{equation}
  \delta_{ka}=Q_a^\top\Gamma_a^{-1/2}(\mu_{ka}-\bar\mu_a),
  \quad R_{ka}=Q_a^\top\widetilde V_{ka}.
\end{equation}
Thus $d_a\leq\sum_kr_k+K-1$.  Numerical SVD truncation is approximate if it
drops a nonzero direction; we log rank and tolerance.

\subsection{Outcome update and reverse diffusion}

For either a fantasy or real $s_a$, component $k$ updates in coefficient
space:
\begin{align}
 \Lambda_{ka}&=(I+R_{ka}^\top R_{ka})^{-1},\\
 \widehat\xi_{ka}
 &=\Lambda_{ka}R_{ka}^\top(s_a-\delta_{ka}),\\
 m_k^+&=m_k+U_k\widehat\xi_{ka},\qquad
 U_k^+=U_k\operatorname{chol}(\Lambda_{ka}), \label{eq:update}\\
 w_k^+&\propto w_k\,
 \N(s_a;\delta_{ka},I+R_{ka}R_{ka}^\top).
\end{align}
The same update changes every unobserved candidate mean through
$V_k\widehat\xi_{ka}$.  After a real observation, $m_k^+$ enters the next
deterministic reverse step.  For VP/DDIM~\cite{song2021ddim},
\begin{equation}
 x_{t'}^k=\alpha_{t'}m_k^+
 +\sigma_{t'}\frac{x_t^k-\alpha_tm_k}{\sigma_t}.
 \label{eq:ddim-update}
\end{equation}
The implementation then re-evaluates the denoiser at the next ordinary step.

\subsection{Receding-horizon planning}

The $H=2$ value of root $a$ is
\begin{equation}
 Q(a)=g_b(a)+\gamma\E_{s_a}
 \left[\max\{0,\max_{a'\in\mathcal A(a)}g_{b^+(s_a)}(a')\}\right],
 \label{eq:h2}
\end{equation}
where $g_b(a)=\I_b(C;y_a)-c_a$.  Projected fantasies use
Eqs.~\eqref{eq:projected-observation}--\eqref{eq:update}, so the tree calls no
denoiser.  The confirmatory configuration supplies and scores every feasible
root, then expands an immediate-utility beam, executes one root, and replans
after its outcome.  The expansion beam is heuristic; any root cap is a
separately labeled engineering ablation.  Immediate bounds safely prune
equal-cost future leaves, never roots, and all-expansion small-scale runs
measure beam sensitivity.

\begin{algorithm}[t]
\caption{\method{} at one acquisition checkpoint}
\label{alg:raosense}
\begin{algorithmic}[1]\footnotesize
\Require states $\{x_t^k,w_k\}$, acquired data $y_M$, feasible groups
\For{$k=1,\ldots,K$}
  \State $m_k\gets D_\theta(x_t^k,t)$; obtain aligned $(U_k,V_k)$
\EndFor
\State build mixture bounds and projected action models
\State score every feasible root; expand $H=2$ beam with projected fantasies
\State acquire only $a^\star$ and observe real $y_{a^\star}$
\State update $\{w_k,m_k,U_k\}$ analytically by Eq.~\eqref{eq:update}
\State continue each branch by Eq.~\eqref{eq:ddim-update}
\end{algorithmic}
\end{algorithm}

\subsection{Target information under operator nuisance}
\label{sec:target-information}

Let $\phi$ be a low-dimensional nuisance in
$y_a=A_a(\phi)x+\epsilon$.  Linearizing around $(m_k,\bar\phi_k)$ gives
\begin{equation}
  y_a\approx A_a(\bar\phi_k)x+
  \left[\partial_\phi A_a(\bar\phi_k)m_k\right](\phi-\bar\phi_k)+\epsilon.
  \label{eq:nuisance-linearization}
\end{equation}
Append a gauge-fixed nuisance factor to each branch.  For one branch, write
the whitened response as
$R_x\xi_x+R_\phi\xi_\phi+e$, initially with independent image and nuisance
coefficients.  Information about the image while marginalizing nuisance is
\begin{equation}
 \I(\xi_x;y_a\mid k)
 =\tfrac12\log\frac{
   \det(I+[R_x,R_\phi]^\top[R_x,R_\phi])}{
   \det(I+R_\phi^\top R_\phi)}.
 \label{eq:target-info}
\end{equation}
After conditioning rotates coefficient directions, we instead use the null
space of the target-state factor.  For a mixture, target $(Z,x)$ retains the
mode term:
\begin{equation}
 \I(Z,x;y_a)
 =\I(Z;y_a)+\sum_kw_k\I(x;y_a\mid Z=k).
 \label{eq:mixture-target}
\end{equation}
This rewards mode/anatomy resolution while treating calibration variation as
structured noise.  Because label information counts, $Z$ must be
target-relevant; nuisance-only or duplicate labels can overcount image
information.  Label-marginalized $\I(x;y_a)$ is not implemented.  This does
not prove global blind identifiability; nuisance bases remain low-dimensional
and gauge-fixed.
